% GENERATED FILE. Edit canonical lessons, then run npm run generate:books.
% canonical-source-hash: fnv1a64:06152af4c1d23bb3
% canonical-lessons: PT-C292-fugir, PT-C292-escapar, PT-C292-fuga, PT-C292-fugitivo, PT-C292-perseguicao, PT-R292-revisao-285-292

\chapter{Fugir da polícia}
\label{ch:pt-fugir-da-policia}
\hlchaptermodality{292}

\begin{chapteropening}
\textbf{By the end of this chapter:} \emph{I can tell how someone ran away and escaped, and talk about the escape and the chase.}

Complete the last lesson of chapter 292: I can tell how someone ran away and escaped, and talk about the escape and the chase.
\end{chapteropening}

\section[fugir]{fugir — to run away, to flee}

\label{lesson:PT-C292-fugir}



\begin{quote}
Before the new one: say the Portuguese for to hold hostage, then the Portuguese for the hostage.
\end{quote}

\subsection*{You'll want to know: fugir}
\textbf{fugir} — \textquotedblleft{}to run away, to flee\textquotedblright{}.

\begin{quote}
O cão fugiu de casa.
\end{quote}

\begin{quote}
\emph{The dog ran away from home.}
\end{quote}

Present: \emph{fujo, foges, foge}.

\begin{cousinweb}
\emph{fugir} is inherited from Vulgar Latin \emph{fugīre}, a remake of Latin \emph{fugere} \textquotedblleft{}to flee\textquotedblright{} (EtymDB).
\end{cousinweb}

\begin{grammarlens}[title={What you've built}]
One more word for telling, explaining and describing at B1.
\end{grammarlens}

\subsection*{Your turn}
\begin{itemize}
\raggedright
  \item \emph{Say it:} \emph{fugir}
  \item \emph{Say it:} \emph{fugir}, once more
  \item \emph{Say it:} a sentence of your own with \emph{fugir}
  \item \emph{Recall:} say the Portuguese for to hold hostage, then the Portuguese for the hostage, then say \emph{fugir} again
\end{itemize}

\begin{tcolorbox}[breakable,colback=teal!4,colframe=teal!35!black,title={Before you move on}]
What does fugir mean? (\textquotedblleft{}To run away, to flee\textquotedblright{}.) Say it once more.
\end{tcolorbox}
\section[escapar]{escapar — to escape, to get away}

\label{lesson:PT-C292-escapar}



\begin{quote}
Before the new one: say the Portuguese for the hostage, then the Portuguese for to run away.
\end{quote}

\subsection*{You'll want to know: escapar}
\textbf{escapar} — \textquotedblleft{}to escape, to get away\textquotedblright{}.

\begin{quote}
O pássaro escapou da gaiola.
\end{quote}

\begin{quote}
\emph{The bird escaped from the cage.}
\end{quote}

\emph{Escapar por pouco} — to have a narrow escape.

\begin{cousinweb}
\emph{escapar} comes down from Old Portuguese (EtymDB lists it with Italian \emph{scappare}); it is usually traced to Vulgar Latin \emph{excappāre} \textquotedblleft{}to slip out of the cape\textquotedblright{}, from \emph{cappa} \textquotedblleft{}a cape\textquotedblright{}.
\end{cousinweb}

\begin{grammarlens}[title={What you've built}]
One more word for telling, explaining and describing at B1.
\end{grammarlens}

\subsection*{Your turn}
\begin{itemize}
\raggedright
  \item \emph{Say it:} \emph{escapar}
  \item \emph{Say it:} \emph{escapar}, once more
  \item \emph{Say it:} a sentence of your own with \emph{escapar}
  \item \emph{Recall:} say the Portuguese for the hostage, then the Portuguese for to run away, then say \emph{escapar} again
\end{itemize}

\begin{tcolorbox}[breakable,colback=teal!4,colframe=teal!35!black,title={Before you move on}]
What does escapar mean? (\textquotedblleft{}To escape, to get away\textquotedblright{}.) Say it once more.
\end{tcolorbox}
\section[a fuga]{a fuga — the escape, the flight}

\label{lesson:PT-C292-fuga}



\begin{quote}
Before the new one: say the Portuguese for to run away, then the Portuguese for to escape.
\end{quote}

\subsection*{You'll want to know: a fuga}
\textbf{a fuga} — \textquotedblleft{}the escape, the flight\textquotedblright{}.

\begin{quote}
A fuga do assaltante durou pouco.
\end{quote}

\begin{quote}
\emph{The robber’s escape did not last long.}
\end{quote}

\emph{Pôr-se em fuga} — to take flight. Also a leak: \emph{uma fuga de gás}.

\begin{cousinweb}
\emph{fuga} comes from Latin \emph{fuga} \textquotedblleft{}flight, escape\textquotedblright{} (EtymDB), from \emph{fugere} \textquotedblleft{}to flee\textquotedblright{}.
\end{cousinweb}

\begin{grammarlens}[title={What you've built}]
One more word for telling, explaining and describing at B1.
\end{grammarlens}

\subsection*{Your turn}
\begin{itemize}
\raggedright
  \item \emph{Say it:} \emph{a fuga}
  \item \emph{Say it:} \emph{a fuga}, once more
  \item \emph{Say it:} a sentence of your own with \emph{a fuga}
  \item \emph{Recall:} say the Portuguese for to run away, then the Portuguese for to escape, then say \emph{a fuga} again
\end{itemize}

\begin{tcolorbox}[breakable,colback=teal!4,colframe=teal!35!black,title={Before you move on}]
What does a fuga mean? (\textquotedblleft{}The escape, the flight\textquotedblright{}.) Say it once more.
\end{tcolorbox}
\section[o fugitivo]{o fugitivo — the fugitive}

\label{lesson:PT-C292-fugitivo}



\begin{quote}
Before the new one: say the Portuguese for to escape, then the Portuguese for the escape.
\end{quote}

\subsection*{You'll want to know: o fugitivo}
\textbf{o fugitivo} — \textquotedblleft{}the fugitive\textquotedblright{}.

\begin{quote}
O fugitivo escondeu-se numa quinta.
\end{quote}

\begin{quote}
\emph{The fugitive hid on a farm.}
\end{quote}

A woman is \emph{a fugitiva}.

\begin{cousinweb}
\emph{fugitivo} was borrowed from Latin \emph{fugitīvus} \textquotedblleft{}runaway\textquotedblright{} (EtymDB), from \emph{fugere} \textquotedblleft{}to flee\textquotedblright{}.
\end{cousinweb}

\begin{grammarlens}[title={What you've built}]
One more word for telling, explaining and describing at B1.
\end{grammarlens}

\subsection*{Your turn}
\begin{itemize}
\raggedright
  \item \emph{Say it:} \emph{o fugitivo}
  \item \emph{Say it:} \emph{o fugitivo}, once more
  \item \emph{Say it:} a sentence of your own with \emph{o fugitivo}
  \item \emph{Recall:} say the Portuguese for to escape, then the Portuguese for the escape, then say \emph{o fugitivo} again
\end{itemize}

\begin{tcolorbox}[breakable,colback=teal!4,colframe=teal!35!black,title={Before you move on}]
What does o fugitivo mean? (\textquotedblleft{}The fugitive\textquotedblright{}.) Say it once more.
\end{tcolorbox}
\section[a perseguição]{a perseguição — the chase, the pursuit}

\label{lesson:PT-C292-perseguicao}



\begin{quote}
Before the new one: say the Portuguese for the escape, then the Portuguese for the fugitive.
\end{quote}

\subsection*{You'll want to know: a perseguição}
\textbf{a perseguição} — \textquotedblleft{}the chase, the pursuit\textquotedblright{}.

\begin{quote}
Houve uma perseguição pelas ruas da cidade.
\end{quote}

\begin{quote}
\emph{There was a chase through the city streets.}
\end{quote}

Also persecution: \emph{a perseguição religiosa}.

\begin{cousinweb}
\emph{perseguição} is built on \emph{perseguir} \textquotedblleft{}to chase\textquotedblright{}, from Latin \emph{persequī} \textquotedblleft{}to follow closely\textquotedblright{}, built on \emph{sequī} \textquotedblleft{}to follow\textquotedblright{} (Lewis \& Short).
\end{cousinweb}

\begin{grammarlens}[title={What you've built}]
One more word for telling, explaining and describing at B1.
\end{grammarlens}

\subsection*{Your turn}
\begin{itemize}
\raggedright
  \item \emph{Say it:} \emph{a perseguição}
  \item \emph{Say it:} \emph{a perseguição}, once more
  \item \emph{Say it:} a sentence of your own with \emph{a perseguição}
  \item \emph{Recall:} say the Portuguese for the escape, then the Portuguese for the fugitive, then say \emph{a perseguição} again
\end{itemize}

\begin{tcolorbox}[breakable,colback=teal!4,colframe=teal!35!black,title={Before you move on}]
What does a perseguição mean? (\textquotedblleft{}The chase, the pursuit\textquotedblright{}.) Say it once more.
\end{tcolorbox}
\section[(dialogue)]{Revisão dos capítulos 285–292}

\label{lesson:PT-R292-revisao-285-292}



\begin{quote}
Say the Portuguese for to interview and the interviewee.
\end{quote}

\subsection*{Your turn}
Say these aloud:

\begin{itemize}
\raggedright
  \item to interview, the interviewee, the listener, the TV viewer, the reader
  \item the offence, the vandal, vandalism, the victim, the homicide
  \item the robbery, the theft, to steal, the pickpocket, the disappearance
  \item the hold-up, to rob, the robber, to break open, the break-in
  \item the fraud, the scam, to swindle, the con man, the accomplice
  \item the weapon, the pistol, to threaten, the threat, the murderer
  \item the kidnapping, to kidnap, the hostage-taking, to hold hostage, the hostage
  \item to run away, to escape, the escape, the fugitive, the chase
\end{itemize}

\begin{tcolorbox}[breakable,colback=teal!4,colframe=teal!35!black,title={Before you move on}]
To interview? (\textbf{entrevistar}.) The interviewee? (\textbf{o entrevistado}.) The listener? (\textbf{o ouvinte}.) The TV viewer? (\textbf{o telespectador}.) The reader? (\textbf{o leitor}.) The offence? (\textbf{o delito}.) The vandal? (\textbf{o vândalo}.) Vandalism? (\textbf{o vandalismo}.) The victim? (\textbf{a vítima}.) The homicide? (\textbf{o homicídio}.) The robbery? (\textbf{o roubo}.) The theft? (\textbf{o furto}.) To steal? (\textbf{furtar}.) The pickpocket? (\textbf{o carteirista}.) The disappearance? (\textbf{o desaparecimento}.) The hold-up? (\textbf{o assalto}.) To rob? (\textbf{assaltar}.) The robber? (\textbf{o assaltante}.) To break open? (\textbf{arrombar}.) The break-in? (\textbf{o arrombamento}.) The fraud? (\textbf{a fraude}.) The scam? (\textbf{a burla}.) To swindle? (\textbf{burlar}.) The con man? (\textbf{o burlão}.) The accomplice? (\textbf{o cúmplice}.) The weapon? (\textbf{a arma}.) The pistol? (\textbf{a pistola}.) To threaten? (\textbf{ameaçar}.) The threat? (\textbf{a ameaça}.) The murderer? (\textbf{o assassino}.) The kidnapping? (\textbf{o rapto}.) To kidnap? (\textbf{raptar}.) The hostage-taking? (\textbf{o sequestro}.) To hold hostage? (\textbf{sequestrar}.) The hostage? (\textbf{o refém}.) To run away? (\textbf{fugir}.) To escape? (\textbf{escapar}.) The escape? (\textbf{a fuga}.) The fugitive? (\textbf{o fugitivo}.) The chase? (\textbf{a perseguição}.)
\end{tcolorbox}
